(1) 一个闭合超导环含两条并联约瑟夫森弱结，临界电流 \(I_1,I_2>0\)，约瑟夫森能量 \(E_j=\hbar I_j/(2e)\)。忽略环自感、温度和结电容，并以 \(\phi=2\pi\Phi_{\rm ext}/\Phi_0\) 定义穿环磁通相位。选 \(\delta_1=\delta+\phi/2,\delta_2=\delta-\phi/2\)，求总超电流 \(I(\delta,\phi)\) 及有效Josephson能量。
(2) 对任意非对称 \(I_1,I_2\)，严格求最大可承载无电压直流电流 \(I_c(\Phi_{\rm ext})\)；分别讨论等结、完全不等结、零磁通和半磁通量子，并证明磁通周期为 \(\Phi_0\)。
(3) 令外电流为 \(I_b\)，写出有效倾斜势 \(U(\delta)\)。当 \(|I_b|<I_c\) 时求单个亚稳极小值和相邻鞍点间的准确逃逸能垒 \(\Delta U\)，给出 \(I_b\to I_c^{-}\) 的 \(3/2\) 幂律，并说明忽略结电容不影响静态势垒。
(4) 在 \(I_b=0\) 时由最低Josephson能量求平衡环流 \(I_{\rm circ}=-\partial F/\partial\Phi_{\rm ext}\)，明确 \(\Phi_{\rm ext}=\Phi_0/2\) 的等结退化点需要单独讨论。检验 \(I_2\to0\) 与 \(I_1=I_2\) 极限，交代自感不可忽略时磁通约束的修正。
